Il capitolo precedente ha definito la parametrizzazione STERA e ha mostrato come una singola proiezione esegue il forward e la backpropagation. Questo capitolo descrive la struttura del modello addestrato e valutato nella tesi, \textbf{STERA 30M}: un Transformer decoder-only in cui tutte le matrici principali sono proiezioni STERA. Il nome indica la classe dimensionale dei \num{30346182} parametri allenabili effettivamente memorizzati, non il numero di elementi delle matrici espanse né una nuova famiglia di architetture. Descriviamo nell'ordine l'architettura, le matrici del modello, il percorso di un token dall'ingresso alla distribuzione sul token successivo, il training e l'esecuzione.

\section{Architettura}
\label{sec:architettura-stera30m}

STERA 30M mantiene la struttura di un Transformer decoder-only e cambia solo la rappresentazione delle sue matrici. Le dimensioni sono:
\begin{itemize}
\item sei layer, con hidden size $d=768$;
\item attention a 16 head da 48 componenti; i layer si alternano fra attention causale sull'intero prefisso e attention causale con finestra locale di 256 posizioni, tre e tre;
\item un feed-forward con dimensione interna $f=\num{3072}$ e activation function GELU;
\item un vocabolario di \num{50257} token, una tabella di \num{2048} posizioni assolute allenabili e finestre di training fino a 512 token.
\end{itemize}

Ogni layer segue lo schema con LayerNorm in ingresso (\emph{pre-norm}): il vettore passa per una LayerNorm, per l'attention e viene sommato alla residual connection; poi passa per una seconda LayerNorm, per il feed-forward e viene di nuovo sommato alla residual connection. Dopo il sesto layer una LayerNorm finale produce la rappresentazione da cui si calcolano i logit del vocabolario. La \autoref{fig:stera30m-flusso}, subito sotto, mostra il percorso completo; in azzurro sono indicate le proiezioni STERA.

\begin{figure}[H]
\centering
\resizebox{.92\linewidth}{!}{\begin{tikzpicture}[
  font=\small, >={Latex[length=2mm]},
  box/.style={draw=steraGrigio, rounded corners=2pt, text width=60mm, minimum height=7mm, align=center, fill=white},
  rm/.style={box, fill=steraBlu!12, draw=steraBlu},
  op/.style={circle, draw=steraGrigio, inner sep=1pt, font=\footnotesize},
  lab/.style={font=\footnotesize, text=steraGrigio, align=left}]
\node[box] (in) {Token $u_t$ e posizione $t$};
\node[rm, above=5mm of in] (emb) {Embedding: $x_t = E_{u_t,:} + P_{t,:}$};
\node[box, above=8mm of emb] (ln1) {LayerNorm};
\node[rm, above=3mm of ln1] (att) {Attenzione causale, 16 teste\\[-1pt]{\footnotesize $W_Q, W_K, W_V, W_O$ RowMask23}};
\node[op, above=3mm of att] (s1) {$+$};
\node[box, above=3mm of s1] (ln2) {LayerNorm};
\node[rm, above=3mm of ln2] (ff) {Feed-forward con GELU, $768\to3072\to768$\\[-1pt]{\footnotesize $W_1, W_2$ RowMask23}};
\node[op, above=3mm of ff] (s2) {$+$};
\node[box, above=8mm of s2] (lnf) {LayerNorm finale};
\node[rm, above=5mm of lnf] (out) {Logit: $z_t = E\,h_t$ \quad {\footnotesize (50.257 candidati)}};
\node[box, above=5mm of out] (p) {$p(u_{t+1}\mid u_{\le t}) = \mathrm{softmax}(z_t)$};
\draw[->] (in) -- (emb); \draw[->] (emb) -- (ln1); \draw[->] (ln1) -- (att); \draw[->] (att) -- (s1);
\draw[->] (s1) -- (ln2); \draw[->] (ln2) -- (ff); \draw[->] (ff) -- (s2); \draw[->] (s2) -- (lnf);
\draw[->] (lnf) -- (out); \draw[->] (out) -- (p);
\coordinate (r0) at ($(emb.north)!.5!(ln1.south)$);
\draw[->] (r0) -- ++(-38mm,0) |- (s1.west);
\coordinate (r1) at ($(s1.north)!.5!(ln2.south)$);
\draw[->] (r1) -- ++(-38mm,0) |- (s2.west);
\begin{scope}[on background layer]
\node[draw=steraGrigio, dashed, rounded corners=4pt, fit=(ln1)(att)(s1)(ln2)(ff)(s2), inner xsep=10mm, inner ysep=3mm] (blk) {};
\end{scope}
\node[lab, anchor=north west] at ([xshift=2mm,yshift=-2mm]blk.north east) {$\times 6$ strati};
\node[lab, anchor=north west, text width=27mm] at ([xshift=2mm,yshift=-8mm]blk.north east) {strati 1, 3, 5: attenzione su tutto il prefisso\\ strati 2, 4, 6: finestra locale di 256 posizioni};
\draw[->, steraBlu, dashed] (emb.west) -- ++(-20mm,0) |- node[pos=.25, left, font=\footnotesize, text=steraBlu]{stessa $E$} (out.west);
\end{tikzpicture}
}
\caption{Architettura di STERA 30M, dal token in ingresso alla distribuzione sul token successivo. In azzurro le matrici rappresentate come proiezioni STERA; la matrice $E$ è condivisa fra embedding e output head.}
\label{fig:stera30m-flusso}
\end{figure}

\section{Le matrici del modello}
\label{sec:atlante-matrici}

STERA 30M ha 37 matrici principali, tutte proiezioni STERA nel senso della \autoref{def:proiezione-stera}, con rango $r=96$ per il termine low-rank:
\begin{itemize}
\item la matrice $E\in\mathbb{R}^{\num{50304}\times768}$, condivisa fra embedding e output head; le righe sono \num{50304} invece di \num{50257} per arrivare a un multiplo di 128, e le 47 righe in più non corrispondono a nessun token;
\item in ognuno dei sei layer, le quattro proiezioni dell'attention $W_Q$, $W_K$, $W_V$, $W_O$, di dimensione $768\times768$;
\item in ognuno dei sei layer, le due proiezioni del feed-forward: $W_1$, di dimensione $\num{3072}\times768$, che espande il vettore, e $W_2$, di dimensione $768\times\num{3072}$, che lo riporta alla dimensione del modello.
\end{itemize}
Ogni matrice è un mosaico di blocchi $128\times128$, come mostra la \autoref{fig:atlante-matrici}, subito sotto; la \autoref{tab:atlante-matrici} conta blocchi, coordinate e valori della forma densa.

\begin{figure}[H]
\centering
\begin{tikzpicture}[x=1.6mm, y=1.6mm, font=\small,
  cel/.style={draw=steraBlu!70, fill=steraBlu!10, line width=.2pt}]
% E: 393 x 6 blocchi, disegnata compressa (24 righe rappresentative)
\foreach \i in {0,...,5} \foreach \j in {0,...,23} \draw[cel] (\i,-\j) rectangle ++(1,-1);
\node[font=\footnotesize] at (3,-25.6) {$\vdots$};
\foreach \i in {0,...,5} \foreach \j in {27,28} \draw[cel] (\i,-\j) rectangle ++(1,-1);
\node[align=center, anchor=north] at (3,-30) {$E$\\{\footnotesize $50.304\times768$}\\{\footnotesize $393\times6$ blocchi}};
% Q, K, V, O: 6x6
\foreach \m/\x in {Q/12, K/21, V/30, O/39} {
  \foreach \i in {0,...,5} \foreach \j in {0,...,5} \draw[cel] (\x+\i,-\j) rectangle ++(1,-1);
  \node[anchor=north] at (\x+3,-7) {$W_{\m}$};
}
\node[anchor=north, font=\footnotesize, align=center] at (28.5,-10.5) {attenzione: $768\times768$, $6\times6$ blocchi ciascuna};
% FF1 24 x 6, FF2 6 x 24
\foreach \i in {0,...,5} \foreach \j in {0,...,23} \draw[cel] (12+\i,-15-\j) rectangle ++(1,-1);
\node[anchor=north, align=center] at (15,-40) {$W_1$\\{\footnotesize $3.072\times768$}};
\foreach \i in {0,...,23} \foreach \j in {0,...,5} \draw[cel] (22+\i,-15-\j) rectangle ++(1,-1);
\node[anchor=west, align=left] at (46.5,-18) {$W_2$\\{\footnotesize $768\times3.072$}};
\draw[decorate, decoration={brace, amplitude=4pt}, steraGrigio] (11.5,1) -- (46.5,1) node[midway, above=4pt, font=\footnotesize] {ripetute in ciascuno dei 6 strati};
\node[anchor=north west, font=\footnotesize, align=left, text width=40mm] at (22,-25) {ogni quadrato è un blocco\\RowMask23 $128\times128$ con\\3.176 coordinate e una scala};
\end{tikzpicture}

\caption{Le matrici di STERA 30M divise in blocchi RowMask23 $128\times128$: la matrice condivisa $E$ e, per ciascuno dei sei layer, le quattro proiezioni dell'attention e le due del feed-forward. Ogni quadrato è un blocco.}
\label{fig:atlante-matrici}
\end{figure}

\begin{table}[H]
\centering
\small
\begin{tabular}{lrrrrr}
\toprule
Matrici & Numero & Dimensione & Blocchi & Coordinate & Valori densi \\
\midrule
$E$ (embedding e output head) & 1 & $\num{50304}\times768$ & \num{2358} & \num{7489008} & \num{38633472} \\
$W_Q,W_K,W_V,W_O$ & 24 & $768\times768$ & \num{864} & \num{2744064} & \num{14155776} \\
$W_1$ & 6 & $\num{3072}\times768$ & \num{864} & \num{2744064} & \num{14155776} \\
$W_2$ & 6 & $768\times\num{3072}$ & \num{864} & \num{2744064} & \num{14155776} \\
\midrule
Totale & 37 & & \num{4950} & \num{15721200} & \num{81100800} \\
\bottomrule
\end{tabular}
\caption{Matrici di STERA 30M: blocchi RowMask23, coordinate (\num{3176} per blocco) e valori della forma densa. Le righe delle proiezioni dei layer sommano i sei layer.}
\label{tab:atlante-matrici}
\end{table}

Alle coordinate si aggiungono le altre componenti delle 37 proiezioni, una scala per ognuno dei \num{4950} blocchi, \num{134016} guadagni e \num{12865536} valori dei fattori $U$ e $V$, e i parametri che restano densi: \num{1572864} posizioni, \num{19968} parametri delle LayerNorm e \num{27648} bias, presenti in $W_O$, $W_1$ e $W_2$. Il totale è di \num{30346182} parametri allenabili, tutti memorizzati in FP32, contro \num{82721280} valori della forma densa equivalente; la composizione per componente e il risparmio del $63{,}32\%$ sono nella \autoref{tab:componenti-stera30m}.

\section{Dal token alla distribuzione sul token successivo}
\label{sec:percorso-token}

Il forward del modello applica le proiezioni STERA nell'ordine dell'architettura. Per una posizione $t$ con token $u_t$:
\begin{enumerate}
\item \textbf{Ingresso.} Si prende la riga $E_{u_t,:}$ della matrice condivisa e le si somma il vettore di posizione $P_{t,:}$, ottenendo un vettore di 768 componenti.
\item \textbf{Layer.} Il vettore attraversa i sei layer. In ciascuno, le proiezioni $W_Q$, $W_K$, $W_V$ producono query, key e value, $W_O$ riporta il risultato dell'attention nello spazio del modello, $W_1$ e $W_2$ formano il feed-forward. Ogni proiezione è applicata con i quattro passi della \autoref{sec:forward-stera}.
\item \textbf{Uscita.} Dopo la LayerNorm finale, la rappresentazione $h_t$ è confrontata con tutte le righe valide di $E$: il logit del token $v$ è $z_{t,v}=E_{v,:}\,h_t$, e la softmax sui \num{50257} logit dà la probabilità del token successivo,
\[
p_{t,v}=\frac{\exp(z_{t,v})}{\sum_{j=0}^{50256}\exp(z_{t,j})}.
\]
\end{enumerate}
La stessa matrice $E$ fornisce quindi sia il vettore d'ingresso di un token sia il suo logit in uscita: è memorizzata una sola volta, ma è usata in due punti del modello.

\section{Training}
\label{sec:addestramento-stera30m}

\paragraph{Loss.} In ogni posizione con token target $y_t$ la loss è $\mathcal L_t=-\log p_{t,y_t}$, mediata sui token target del batch ed escluso il padding.

\paragraph{Backpropagation nel modello.} Il gradiente parte dai logit: rispetto al logit del token $v$ vale $\partial\mathcal L_t/\partial z_{t,v}=p_{t,v}-\mathbf 1\{v=y_t\}$, cioè la probabilità assegnata meno uno per il token corretto. Da qui risale nell'ordine inverso del forward: attraversa la output head, la LayerNorm finale e i sei layer, dall'ultimo al primo. Ogni proiezione STERA che incontra riceve il gradiente della propria uscita, $\delta=\partial\mathcal L/\partial y$, e con i passi B1--B4 della \autoref{sec:backward-passi} produce i gradienti delle proprie coordinate, scale, guadagni e fattori $U$ e $V$, e il gradiente del proprio ingresso, che prosegue verso il layer precedente. La matrice condivisa $E$ riceve due contributi, che si sommano: quello della output head, $\sum_t(p_t-\mathbf 1_{y_t})\,h_t^\top$ sulle righe valide, e quello delle righe usate come ingresso, che arriva in fondo al percorso.

\subsection{Un passo di training: che cosa viene aggiornato}
\label{sec:passo-training}

Durante il training la matrice densa $W$ compare, ma solo come strumento di calcolo: i parametri che il modello impara e che vengono salvati sono le coordinate compresse, insieme a scale, guadagni e fattori $U$ e $V$. Un passo di training si svolge in cinque fasi, mostrate nella \autoref{fig:passo-training}, subito sotto.

\begin{figure}[H]
\centering
\resizebox{\linewidth}{!}{\begin{tikzpicture}[x=10mm, y=10mm, font=\small, >={Latex[length=2mm]},
  par/.style={draw=steraBlu, fill=steraBlu!12, rounded corners=3pt, minimum width=43mm, minimum height=14mm, align=center},
  tmp/.style={draw=steraGrigio, dashed, fill=steraGrigio!8, rounded corners=3pt, minimum width=43mm, minimum height=14mm, align=center},
  gr/.style={draw=steraArancio, fill=steraArancio!12, rounded corners=3pt, minimum width=43mm, minimum height=14mm, align=center},
  gt/.style={draw=steraArancio, dashed, fill=steraArancio!8, rounded corners=3pt, minimum width=43mm, minimum height=14mm, align=center},
  et/.style={font=\footnotesize, align=center, fill=white, inner sep=1.5pt}]
\node[par] (p) at (0,0) {\textbf{parametri salvati}\\coordinate $c$, scale $\alpha$,\\guadagni $g$, $U$, $V$};
\node[tmp] (w) at (7.6,0) {\textbf{matrice $W$}\\densa, temporanea\\\num{16384} valori per blocco};
\node[tmp] (l) at (15.2,0) {\textbf{loss} $\mathcal L$\\calcolata con $W$};
\node[gt] (g) at (15.2,-3.6) {\textbf{gradiente $G$}\\uno per casella\\\num{16384} per blocco};
\node[gr] (gc) at (6.6,-3.6) {\textbf{gradiente dei parametri}\\uno per coordinata\\\num{3176} per blocco};
\draw[->, steraBlu, thick] (p) -- node[et, above=2pt] {1. ricostruzione} (w);
\draw[->, steraBlu, thick] (w) -- node[et, above=2pt] {2. forward} (l);
\draw[->, steraArancio, thick] (l) -- node[et, right=2pt] {3. backward} (g);
\draw[->, steraArancio, thick] (g) -- node[et, above=2pt] {4. somma per coordinata} (gc);
\draw[->, steraArancio, thick] (gc) -| node[et, pos=.3, above=2pt] {5. AdamW} (p);
\end{tikzpicture}
}
\caption{Un passo di training di STERA 30M. In blu il forward, in arancione il backward. Le caselle tratteggiate sono temporanee: la matrice densa $W$ e il suo gradiente $G$ vengono ricalcolati a ogni passo e non sono mai salvati né aggiornati.}
\label{fig:passo-training}
\end{figure}

\begin{enumerate}
\item \textbf{Ricostruzione.} Dalle coordinate, dalle scale, dai guadagni e dai fattori $U$ e $V$ si costruisce la matrice densa $W$ di ogni proiezione, secondo la \autoref{def:proiezione-stera}. In un blocco, la coordinata di base $b_8$ viene copiata nelle 128 caselle con $i\oplus j=8$.
\item \textbf{Forward.} Il modello usa le matrici $W$ così costruite per calcolare le probabilità dei token e la loss $\mathcal L$. La loss è quindi calcolata con la forma densa equivalente, che però è determinata interamente dalle coordinate.
\item \textbf{Backward.} L'autograd calcola il gradiente della loss rispetto a ogni casella di $W$: è la matrice $G=\partial\mathcal L/\partial W$, con un valore per ciascuna delle \num{16384} caselle di un blocco. Le 128 caselle di $b_8$ ricevono in generale 128 gradienti diversi.
\item \textbf{Raccolta dei gradienti.} Il gradiente di ogni coordinata è la somma dei gradienti delle caselle in cui la coordinata compare (\autoref{sec:gradiente-coordinate}), ciascuno moltiplicato per i fattori che nel forward moltiplicavano la coordinata in quella casella: la scala del blocco e i guadagni della riga e della colonna. Per $b_8$ è la somma dei 128 contributi delle sue caselle; per una correzione o un elemento della diagonale è il gradiente della sua unica casella. Si ottiene un solo gradiente per coordinata, \num{3176} per blocco invece di \num{16384}. Allo stesso modo scale, guadagni, $U$ e $V$ ricavano il proprio gradiente da $G$, con i passi della \autoref{sec:backward-passi}.
\item \textbf{Aggiornamento.} AdamW aggiorna le coordinate, le scale, i guadagni, $U$, $V$ e i parametri densi del modello (posizioni, LayerNorm e bias). La matrice $W$ e il gradiente $G$ vengono scartati, e al passo successivo $W$ si ricostruisce dalle coordinate aggiornate: le 128 caselle di $b_8$ ricevono di nuovo tutte lo stesso valore, quello nuovo.
\end{enumerate}

La \autoref{tab:raccolta-gradienti}, subito sotto, riassume da quali caselle di un blocco ogni tipo di coordinata raccoglie il proprio gradiente. Una casella che contiene base e correzione, $b_{i\oplus j}+e_{i,k}$, contribuisce con lo stesso valore di $G$ a entrambe le coordinate.

\begin{table}[H]
\centering
\small
\renewcommand{\arraystretch}{1.25}
\begin{tabular}{lll}
\toprule
Coordinate & Caselle da cui raccoglie il gradiente & Gradiente \\
\midrule
\rowcolor{steraBlu!12} diagonale $d_i$ & una: $(i,i)$ & $G_{ii}$ \\
\rowcolor{steraGrigio!12} base $b_s$ & 128: tutte le $(i,\,i\oplus s)$ & $\sum_i G_{i,\,i\oplus s}$ \\
\rowcolor{steraArancio!15} correzione $e_{i,k}$ & una: $(i,\,i\oplus\sigma(i,k))$ & $G_{i,\,i\oplus\sigma(i,k)}$ \\
\bottomrule
\end{tabular}
\caption{Raccolta dei gradienti in un blocco RowMask23: per ogni tipo di coordinata, le caselle della matrice ricostruita i cui gradienti si sommano. $G$ è qui il gradiente rispetto al blocco, che include già scala e guadagni (\autoref{sec:backward-passi}).}
\label{tab:raccolta-gradienti}
\end{table}

\paragraph{Perché si aggiornano le coordinate e non la matrice densa.} Si sommano i gradienti delle caselle e poi si aggiorna una sola volta la coordinata; non si aggiornano le caselle una per una. La differenza è sostanziale. Se l'optimizer aggiornasse direttamente le \num{16384} caselle, ciascuna con il proprio gradiente, le 128 caselle di $b_8$ prenderebbero valori diversi fra loro: la matrice uscirebbe dalla famiglia RowMask23 e non si potrebbe più rappresentare in modo esatto con le coordinate. Aggiornando le coordinate, invece, ogni matrice resta per costruzione una proiezione STERA a ogni passo. Inoltre AdamW non è lineare, perché divide ogni aggiornamento per una stima della varianza del gradiente: sommare i gradienti e poi aggiornare la coordinata non equivale a sommare gli aggiornamenti delle singole caselle.

\paragraph{Implementazione e optimizer.} Il training usa PyTorch su due GPU NVIDIA P100. Le fasi 3 e 4 non sono implementate a mano: la ricostruzione della fase 1 è scritta con operazioni di indicizzazione, e l'autograd, risalendo quelle operazioni, somma automaticamente i gradienti delle caselle che condividono una coordinata. Il risultato coincide con quello della \autoref{sec:backward-stera} (\autoref{prop:backward-esatto}). Il weight decay si applica a coordinate, fattori $U$ e $V$ e posizioni, non a scale, guadagni, LayerNorm e bias. Dati, iperparametri e andamento del training sono descritti nel \autoref{cap:capitoli/05-metodologia}.

\section{Dall'training alla generazione}
\label{sec:guida-inferenza-stera}

Al termine del training i parametri sono fissati, e le stesse proiezioni si possono eseguire in tre modi, mostrati nella \autoref{fig:percorso-inferenza-stera}, subito sotto:
\begin{enumerate}
\item \textbf{espansione residente}: si ricostruiscono una volta tutte le matrici e si tengono in memoria; i prodotti sono quelli densi, ma la memoria occupata è quella della forma densa;
\item \textbf{espansione al bisogno}: si ricostruiscono le matrici di un layer prima di usarle e si rilasciano subito dopo; la memoria resta bassa, ma la ricostruzione si ripete;
\item \textbf{prodotto diretto sulle coordinate}: si applicano i quattro passi della \autoref{sec:forward-stera} senza costruire nessuna matrice.
\end{enumerate}

\begin{figure}[H]
\centering
\begin{tikzpicture}[font=\small, >={Latex[length=2mm]},
  col/.style={draw=steraGrigio, rounded corners=3pt, minimum width=44mm, minimum height=50mm, anchor=north},
  co/.style={draw=steraBlu, fill=steraBlu!15, minimum width=10mm, minimum height=10mm},
  de/.style={draw=steraGrigio, fill=steraGrigio!15, minimum width=22mm, minimum height=22mm},
  tit/.style={font=\bfseries\small, align=center}]
\foreach \x/\t in {0/{Espansione\\residente}, 46/{Espansione\\al bisogno}, 92/{Prodotto diretto\\sulle coordinate}} {
  \node[col] at (\x mm,0) {};
  \node[tit, anchor=north] at (\x mm,-2mm) {\t};
}
% 1
\node[de] at (0,-24mm) {$W$};
\node[font=\footnotesize, align=center] at (0,-42mm) {in memoria: matrice densa\\calcolo: $y=Wx$};
% shift group 2
\begin{scope}[xshift=46mm]
\node[co] (c2b) at (-13mm,-20mm) {$c$};
\node[de, minimum width=14mm, minimum height=14mm, dashed] (w2b) at (12mm,-20mm) {$W$};
\draw[->] (c2b) -- node[above, font=\scriptsize]{espandi} (w2b);
\node[font=\footnotesize, align=center] at (0,-38mm) {in memoria: coordinate\\e una matrice temporanea\\calcolo: espansione $+$ $Wx$};
\end{scope}
\begin{scope}[xshift=92mm]
\node[co] (c3) at (0,-18mm) {$c$};
\node[font=\scriptsize, align=center] at (0,-30mm) {$y = d\odot x + \tfrac{1}{N}H\big((Hb)\odot(Hx)\big)$\\$+\ \text{correzioni} + U(V^{\top}x)$};
\node[font=\footnotesize, align=center] at (0,-42mm) {in memoria: solo coordinate\\calcolo: trasformate veloci};
\end{scope}
\draw[->, steraGrigio] (-21mm,-53mm) -- node[below, font=\footnotesize]{memoria occupata dai pesi decrescente} (113mm,-53mm);
\end{tikzpicture}

\caption{Tre modi di eseguire le stesse proiezioni STERA in inferenza. Da sinistra a destra diminuisce la memoria occupata dai pesi e cresce il lavoro di calcolo per ricostruire o applicare la matrice.}
\label{fig:percorso-inferenza-stera}
\end{figure}

La scelta dipende dalla fase della generazione. Nel \emph{prefill} il modello elabora tutti i token del prompt insieme e costruisce la KV cache: molti token usano gli stessi pesi, il costo è dominato dal calcolo e conviene l'espansione al bisogno, che permette prodotti densi fra matrici. Nel \emph{decode} il modello genera un token alla volta: ogni token rilegge tutti i pesi, il costo è dominato dai byte letti e conviene il prodotto diretto sulle coordinate. Le strategie si possono quindi combinare. In tutti i casi si esegue la stessa funzione con gli stessi parametri appresi: il confronto fra il motore compatto e quello denso del \autoref{cap:capitoli/07-sistemi} non confronta due modelli, ma due modi di eseguire lo stesso modello.

\section{Sintesi}

STERA 30M è un Transformer decoder-only a sei layer con 37 matrici principali, tutte proiezioni STERA, per \num{4950} blocchi RowMask23 e \num{30346182} parametri allenabili. Il forward e la backpropagation del modello attraversano ogni proiezione con i passi della \autoref{sec:backward-stera}; in training le matrici sono ricostruite e i gradienti calcolati con l'autograd, in inferenza le stesse proiezioni si eseguono sulle matrici espanse o direttamente sulle coordinate. Il capitolo successivo descrive i dati, il training e il protocollo con cui il modello è valutato.
